\lecture{Matlab}{02/10/2025}
\section{Matlab Basics}
\gls{mat} è un ambiente di programmazione e calcolo numerico. 

È progettato per:
\begin{itemize}
    \item  manipolazione di matrici e vettori,
    \item  simulazione numerica,
    \item  visualizzazione grafica,
    \item  sviluppo di algoritmi scientifici o ingegneristici.
\end{itemize}
È praticamente una via di mezzo tra un linguaggio di programmazione e una potente calcolatrice scientifica.

Nel testo vengono utilizzati i seguenti file su MATLAB\_R2025b:
\begin{itemize}
    \item  \verb|helloworld.m|
    \item  \verb|matrix.m|
\end{itemize}
%
\subsection{Printing}
Per stampare valori di ritorno, contenuti di variabili etc \dots non si inserisce il \verb|;| al termine di un 
espressione. Non inserendo tale carattere farà stampare in console.
\subsubsection{Printing functions}
Si può generare l'output su schermo tramite diverse funzioni, principalmente tramite: \verb|disp()| e \verb|fprintf()|.
I due metodi differiscono per flessibilità e formato dell'output.
\paragraph{Funzione disp()}
La funzione \verb|disp()| (display) è il metodo più semplice per visualizzare variabili o
messaggi testuali. Non consente la formattazione avanzata ma è utile per una stampa rapida.

\verb|disp('Hello World')| stamperà una stringa, \verb|disp(A)| stampa cosa è contenuto nella variabile A.

Con \verb|disp()| si può stampare unicamente una variabile/stringa, non è possibile quindi concatenare variabili
o stringhe con variabili.   

\begin{note}
Si può osservare che la funzione \verb|disp()| fa veramente cagare.
\end{note}

\paragraph{Funzione fprintf()}
La funzione \verb|fprintf()| (formatted print) permette un controllo più preciso sul formato dell'output.
Può essere utilizzato per stampare testo, numeri e variabili all'interno della stessa stringa.
\begin{verbatim}
fprintf('Il valore di x è: %d\n', x)
fprintf('Il risultato è: %.2f\n', y)
\end{verbatim} 
Dove:
\begin{itemize}
	\item \texttt{\%d} indica un numero intero;
	\item \texttt{\%.2f} indica un numero reale con due cifre decimali;
	\item \texttt{\textbackslash n} va a capo (new line).
\end{itemize}

%
\subsection{Variables}
Non vi è il bisogno di esplicitare il tipo di variabile. Per assegnare un valore ad una variabile: 
\[ 
  x=3 \cdot 2
\] 
In questo caso, la varibile $ x $, autodichiarata come \emph{double} (più a riguardo dopo), avrà conservato all'intero
$ 6 $.

Se non vengono specificate variabili, ad esempio nel caso in cui l'espressione sia solo $ 3 \cdot 2 $, nella console verrà
stampata una variabile direttamente assegnata da Matlab, di nome \emph{``ans''} , che verrà sovrascritta
ogni volta che verrà stampato un valore non contenuto in una variabile. 

\subsubsection{Data types}
In Matlab si utilizzando diversi tipi di variabile, per scopi diversi e con diverse grandezze.
Tra le più utilizzate:
\begin{itemize}
    \item \textbf{double}: Array a doppia precisione
    \item \textbf{single}: Array a precisione singola
    \item \textbf{int8}\dots \textbf{int64}:: Array di numeri interi con segno da 8 fino a 64 bit
    \item \textbf{uint8}\dots \textbf{uint64}:: Array di numeri interi senza segno da 8 fino a 64 bit
\end{itemize}

\marginnote{continua\dots}

%
\subsection{Vectors}
Vi sono due tipi di vettori:
\paragraph{Row vectors} 
Per definire un row vector si utilizza:
\begin{verbatim}
x=[1,2,3] oppure x=[1 2 3]
\end{verbatim}
Entrambe le diciture sono valide, il vettore riga può essere dichiarato sia con le virgole che con gli spazi.

\paragraph{Column vectors}
Per definire un column vector si utilizza:
\begin{verbatim}
x=[1;2;3]
\end{verbatim}
Il vettore colonna, a differenza di quello a riga, può essere definito solamente con i \verb|;|.

\subsubsection{Functions and vector manipulation}
\paragraph{length(x)} 
La funzione \verb|length(x)| ritorna la lunghezza del vettore \verb|x|.
\paragraph{Accessing the k-th value of a vector}
Per accedere al valore k-esimo di un vettore si utilizza:
\begin{verbatim}
x(k), x(end)
\end{verbatim}
Tramite \verb|x(k)| si accede al k-esimo valore del vettore (cioè il valore di indice k, indici contati
a partire da 0), tramite \verb|x(end)| si accede al valore di ultima posizione, ovvero quello con indice più
grande.

\paragraph{Vector manipulation}
Due vettori possono essere concatenati insieme, ammesso che abbiano la stessa dimensione, tramite:
\begin{verbatim}
z[xy], z[x;y]
\end{verbatim} 
Entrambe le diciture sono valide.

Si può eliminare un valore contenuto nel vettore di una certa posizione k tramite:
\begin{verbatim}
x(k)=[]
\end{verbatim} 

\paragraph{Creation of n-dimensional vectors}
Si possono creare vettori n-dimensionali tramite:
\begin{verbatim}
zeros(n,1), zeros(1,n);
ones(n,1), ones(1,n);
\end{verbatim}
Questi creano un vettore di lunghezza e altezza n pieni di 0 (zeros) o 1 (ones). 
Deve essere presente minimo 1 colonna/riga.
\begin{matlab}
zeros(2,5)
ones(5,2)
\end{matlab}
Risulta in output:
\begin{matlab}
ans =
    0   0   0   0   0
    0   0   0   0   0

ans =
    1   1
    1   1
    1   1
    1   1
    1   1
\end{matlab}

Si può creare un vettore riga di n punti equidistanti compresi tra \emph{a} e \emph{b} tramite:
\begin{verbatim}
linspace(a,b,n)
\end{verbatim} 
\begin{example}
    Creazione di un vettore tramite linspace di arg: \verb|1,2,4|.
\end{example} 
\begin{matlab}
linspace(1,2,4)
\end{matlab}
Da come output:
\begin{matlab}
ans =
1.0000  1.3333  1.6667  2.0000
\end{matlab}

\bigskip
Mentre il seguente:
\begin{matlab}
a=1;
b=5;
k=b;
linspace(a,b,k)
\end{matlab}
Da come output:
\begin{matlab}
ans =
1   2   3   4   5
\end{matlab}

\bigskip
Per creare un vettore riga di n punti equidistanti con salti s, si utilizza:
\begin{verbatim}
a:s:b
\end{verbatim} 
dove \verb|s| è il nostro \emph{``step''}. 
\begin{example}
    Creazione di un vettore riga con salti di 0.25. 
\end{example} 
\begin{matlab}
0:0.25:1
\end{matlab}
\begin{matlabconsole}
ans =
0   0.2500  0.5000  0.7500  1.0000
\end{matlabconsole}

\subparagraph{Calculate a vector's p-norm}
Dato un vettore $ x $, si può calcolare la trasposta $ x' $ e la sua p-norm tramite:
\begin{verbatim}
norm(x,p)
\end{verbatim} 
\begin{example}
\end{example} 
\begin{matlab}
% Create vector
k=4;
a = linspace (1,5,k);
b = linspace (3,8,k);
c = linspace (21,34,k);
x =[a;b;c]

% Calc p-norm
norm(x,1)
\end{matlab}
\begin{matlabconsole}
x =
1.0000  2.3333  3.6667  5.0000
3.0000  4.6667  6.3333  8.0000
21.0000 25.3333 29.6667 34.0000

ans =
47
\end{matlabconsole}

\paragraph{Maximun \& minimun}
Il massimo e minimo di un vettore può essere computato tramite:
\begin{verbatim}
max(x), min(x)
\end{verbatim}

\begin{example}
Calcolo del max e del min del vettore $ m $
\end{example} 
\begin{matlab}
y=[1,3,5];
x=[2,4,6];
m=[y;x];

m=max(m)
n=min(m)
\end{matlab}
\begin{matlabconsole}
m =

     2     4     6


n =

     2
\end{matlabconsole}

\subsubsection{Operations with vectors}
Dati due vettori (row vectors) $ x $ ed $ y $ di dimensione $ n $, possiamo computare la loro
somma $ x+y $ ed il loro prodotto scalare $ x \cdot y' $. 
\begin{example}
Calcolo della somma e del prodotto scalare dei vettori $ x $ e $ y $.
\end{example}

\begin{matlab}
x=[1,5,8];
y=[3,6,1];

s=x+y
p=x*y'
\end{matlab}
\begin{matlabconsole}
s =

     4    11     9


p =

    41
\end{matlabconsole}
Come si può vedere, per calcolare il prodotto scalare $ x \cdot y' $ abbiamo dovuto indicare $ y' $ 
nella riga di codice. Si potrebbe utilizzare anche il seguente metodo:
\begin{verbatim}
p = dot(x,y)
\end{verbatim} 
che restituisce lo stesso risultato. L'ultimo metodo menzionato è utile ricordarlo, poiché lo si può
utilizzare anche insieme ad altri tipi di vettori (non solo riga). A riguardo più avanti.

\paragraph{Entry-wise operations}
In Matlab, vi sono due tipi di operazioni fattibili tra matrici:
\begin{enumerate}
    \item Operazioni matriciali: \verb|a*b|
    \item Operazioni entry-wise: \verb|a.*b|
\end{enumerate}
\subparagraph{Matrix operations}
Avendo due matrici $ a $ e $ b $, di grandezza uguale, si possono effettuare operazioni
come moltiplicazione, frazione e potenza che rispettano le regole dell'algebra lineare.
\begin{matlab}
A = [1 2; 3 4];
B = [5 6; 7 8];
A*B          % prodotto tra matrici
\end{matlab}

\subparagraph{Entry-wise operations}
Avendo due matrici $ a $ e $ b $, si possono fare operazioni su singoli elementi corrispondenti (quindi
elemento per elemento) posizionando un \verb|.| prima del simbolo dell'operatore, ad esempio:
\begin{verbatim}
a.*b, a./b, a.^b
\end{verbatim} 
\begin{matlab}
A = [1 2; 3 4];
B = [5 6; 7 8];
A.*B          % Moltiplicazione
\end{matlab}

\section{Matrices}
Per definire una matrice $ A $ si utilizza:
\begin{verbatim}
A=[1 2 3; 4 5 6]
\end{verbatim} 
Scritta in termini matematici:
\[ A = 
  \begin{bmatrix}
  1 & 2 & 3 \\
  4 & 5 & 6
  \end{bmatrix}
\] 

\paragraph{Matrix size}
Per avere la grandezza della matrice A, si utilizza la funzione \verb|size(A)|

\subsection{Special matrices}
Vi sono inoltre le seguenti matrici speciali:
\begin{verbatim}
ones(n,m), zeros(n,m), eye(n)
\end{verbatim} 

\begin{example}
Stampa di matrici speciali (uno, zero e matrice identità).
\end{example} 
\begin{matlab}
% Matrice di soli 1 (3x3)
A = ones(3, 3);
% Matrice di soli 0 (2x4)
B = zeros(2, 4);
% Matrice identita' (4x4)
C = eye(4);

% Stampa risultati
disp('A=ones(3,3):')
disp(A)
disp('B=zeros(2,4):')
disp(B)
disp('C=eye(4):')
disp(C)
\end{matlab}
\begin{matlabconsole}
>> matrix
A = ones(3,3):
     1     1     1
     1     1     1
     1     1     1

B = zeros(2,4):
     0     0     0     0
     0     0     0     0

C = eye(4):
     1     0     0     0
     0     1     0     0
     0     0     1     0
     0     0     0     1
\end{matlabconsole}

\subsection{Functions and matrix manipulation}
Tutte le operazioni già definite per i vettori possono essere utilizzate con le matrici, anche i comandi
utilizzati per la manipolazione, per quest'ultimi bisogna però prestare particolare attenzione alle dimensioni!

\paragraph{Accessing the k-th value of a matrix}
Per accedere al valore k-esimo di una matrice si utilizza:
\begin{verbatim}
A(:.k)
\end{verbatim}
Tramite \verb|A(:.k)| si accede al k-esimo valore della matrice (cioè il valore di indice k, indici contati
a partire da 0).

\paragraph{diag(A)} Avendo una matrice $ A $, la funzione \verb|diag(A)| estrae la diagonale della matrice.
Avendo un vettore $ x $, la funzione \verb|diag(x)| ritorna una matrice diagonale del vettore $ x $.

\paragraph{Eigenvalues and eigenvectors} Prima di tutto, diamo una rinfrescata sulla definizione degli autovalori
e delle matrici inverse:
\begin{definition}[Autovalore e autovettore di una matrice quadrata]
Sia $ A $ una matrice quadra $ n \times n $. 

Un numero $ \lambda \in \R $ e un vettore non nullo $ v \in \R^n $ soddisfano:
\[ 
  Av=\lambda v
\] 

Allora:
\begin{itemize}
    \item $ \lambda $ è un autovalore (eigenvalue) di $ A $;
    \item $ v $ è un autovettore (eigenvector) associato a $ \lambda $.
\end{itemize}
\end{definition}
Dal punto di vista geometrico:
\begin{itemize}
    \item L'autovettore indica una direzione invariata dalla trasformazione.
    \item L'autovalore indica di quanto quella direzione viene scalata:
        \begin{itemize}
            \item se $ \lambda > 1 $: la direzione si allunga;
            \item se $ 0 < \lambda < 1 $: si accorcia;
            \item se $ \lambda < 0 $: si inverte. 
        \end{itemize}
\end{itemize}
\begin{example}
Data la matrice 
\[ 
  A=
  \begin{bmatrix}
  2 & 0 \\
  0 & 3
  \end{bmatrix}
\] calcolare l'autovalore e l'autovettore.
\end{example} 
\begin{matlab}
A = [2 0; 0 3];
eig(A)
inv(A)
\end{matlab}
\begin{matlabconsole}
ans =

     2
     3


ans =

    0.5000         0
         0    0.3333
\end{matlabconsole}

\paragraph{linsolve(A,b)} Dato un sistema lineare $ Ax=b $, dove $ A $ è la matrice dei coefficienti, $ x $
è il vettore incognito e $ b $ è il vettore dei termini noti. 
Il comando \verb|linsolve(A,b)| permette di risolvere il sistema lineare, quindi di trovare l'incognita $ x $.
\begin{example}
Dato il seguente sistema lineare:
\[ 
  \begin{cases}
        2x_1 + x_2 = 5 \\
        3x_1 + 4x_2 = 6
    \end{cases}
\] 
trovare il valore di $ x_{1} $ e $ x_{2} $.
\end{example} 
Sistema lineare in forma matriciale:
\[ A=
    \begin{bmatrix}
        2 & 1 \\
        3 & 4
    \end{bmatrix},
    b=
    \begin{bmatrix}
    5 \\ 6
    \end{bmatrix}
\] 
\begin{matlab}
A = [2 1; 3 4];   % matrice 2x2
b = [5; 6];       % vettore colonna

x = linsolve(A, b)
\end{matlab}
\begin{matlabconsole}
>> matrix

x =

    2.8000
   -0.6000
\end{matlabconsole}
Significa quindi che:
\begin{align*}
    x_{1}&=2.8 \\
    x_{2}&=-0.6
\end{align*}